% !TEX root = ../main.tex
%----------------------------------------------------------
% 03-possibility-primer.tex — Section 2: Possibility Theory Primer
% MIGRATED from preprint 04-possibility-theory.tex
% Ozone categories → SPC categories (MRGL, SLGT, ENH, MDT, HIGH)
% Removed: Clyfar-specific FIS/fuzzification framing, figures
%----------------------------------------------------------

\section{Possibility Theory Primer}\label{sec:poss_primer}

% [PLACEHOLDER: Opening paragraph — frame this section as a self-contained
%  primer for atmospheric scientists unfamiliar with possibility theory.
%  No prior knowledge of fuzzy logic is assumed.  Running example: a
%  possibilistic convective outlook assigning possibility values to the
%  five SPC categories \spc{MRGL}, \spc{SLGT}, \spc{ENH}, \spc{MDT},
%  \spc{HIGH}.]

\subsection{Axioms and Definitions}\label{sec:poss_axioms}

Let $\Omega$ denote the \textit{universe of discourse}---the complete set
of mutually exclusive and collectively exhaustive outcomes for a variable
of interest (e.g., SPC convective outlook categories
$\Omega = \{\spc{MRGL},\, \spc{SLGT},\, \spc{ENH},\, \spc{MDT},\,
\spc{HIGH}\}$).  A \textit{possibility distribution}
$\pi: \Omega \to [0,1]$ assigns to each outcome $\omega \in \Omega$ a
degree of possibility $\pi(\omega)$ satisfying:

\begin{axiomlist}
\item \textbf{Something must be possible:} There exists at least one
    outcome $\omega^* \in \Omega$ such that $\pi(\omega^*) > 0$.
\item \textbf{Normalisation (optional):} If
    $\max_{\omega \in \Omega} \pi(\omega) = 1$, the distribution is
    \textit{normal}; otherwise it is \textit{subnormal}.  Classical
    possibility theory \citep{Dubois1988-nh} requires normalisation: at
    least one outcome must be fully possible for necessity to be
    well-defined.  This paper relaxes that requirement, retaining
    subnormal distributions when the forecasting system provides weak
    support for all outcomes (Section~\ref{sec:subnormal}).
\end{axiomlist}

From a possibility distribution $\pi$ that may or may not be normalised,
two dual measures are defined for any event $A \subseteq \Omega$:

\begin{align}
\Pi(A) &= \max_{\omega \in A} \pi(\omega)
    \quad \text{(Possibility)} \label{eq:poss} \\
N(A) &= 1 - \Pi(\neg A)
     = 1 - \max_{\omega \notin A} \pi(\omega)
    \quad \text{(Necessity)} \label{eq:nec}
\end{align}

\textbf{Interpretation:}
\begin{itemize}[leftmargin=2em]
\item $\Pi(A)=1$: Event $A$ is entirely plausible given current knowledge.
\item $\Pi(A)=0$: Event $A$ is impossible (inconsistent with available
    evidence).
\item $N(A)=1$: Event $A$ is certain (all alternatives are ruled out).
\item $N(A)=0$: Event $A$ is not certain (at least one alternative
    remains plausible).
\end{itemize}

For any event $A$, $N(A) \leq \Pi(A)$ always holds
\citep{Dubois1988-nh}.  The interval $[N(A),\, \Pi(A)]$ brackets the
range of probabilities $P(A)$ consistent with the available information:
any probability distribution compatible with $\pi$ must satisfy
$N(A) \leq P(A) \leq \Pi(A)$.
% [PLACEHOLDER: Add Zadeh and Dubois 2006 citations for the bracketing
%  property.]

%----------------------------------------------------------
\subsection{Subnormal Distributions and Ignorance}\label{sec:subnormal}

% [PLACEHOLDER: Opening paragraph — explain why subnormality matters in
%  severe-weather forecasting.  When a forecasting system has weak support
%  for all SPC categories, the resulting distribution should remain
%  subnormal rather than being normalised to hide this lack of knowledge.]

Define the (possibilistic) \textit{ignorance} measure \ign{} as the
residual possibility ``mass'' not covered in $\pi$:
\begin{equation}
H_\Pi = 1 - \max_{\omega \in \Omega} \pi(\omega) \label{eq:ignorance}
\end{equation}

\textbf{Interpretation:} \ign{} quantifies the fraction of plausibility
unaccounted for by the model's knowledge.  High ignorance
($H_\Pi \approx 1$) signals that the forecasting system encounters
conditions outside its coverage or that conflicting evidence produces
weak support for all outcomes.  Low ignorance ($H_\Pi \approx 0$)
indicates the system assigns high possibility to at least one
outcome---an impossible event ($\Pi(\overline\omega) = 0$) can be ruled
out due to known plausibility of a different outcome, rather than
spurious confidence that stems from enforced normalisation.

% [PLACEHOLDER: Paragraph connecting ignorance to the normalisation
%  problem.  Normalising forces derived probability distributions to sum
%  to unity, masking the distinction between "the system predicts
%  \spc{MDT}" and "the system has insufficient information to rule out any
%  category."  Preserving subnormality trades formal completeness for
%  operational transparency.]

\textbf{Example 1 (sharp \spc{MDT} forecast):}
% [PLACEHOLDER: Describe the meteorological scenario — e.g., a classic
%  severe-weather setup with strong instability, shear, and moisture.]
The resulting possibility distribution is:
\[
\pi = \bigl\{\spc{MRGL}: 0.0,\; \spc{SLGT}: 0.1,\;
    \spc{ENH}: 0.2,\; \spc{MDT}: 0.75,\; \spc{HIGH}: 0.15\bigr\}
\]
with $\ign = 1 - 0.75 = 0.25$ and
$\Nc(\spc{MDT}) = 1 - 0.2/0.75 = 0.73$.  Interpretation: \spc{MDT}
is the most plausible outcome ($\Pi = 0.75$) with moderate ignorance
($\ign = 0.25$); conditional necessity of $0.73$ indicates that among
scenarios the system covers, \spc{MDT} is the dominant category.

\textbf{Example 2 (ambiguous \spc{MRGL}/\spc{SLGT} forecast):}
% [PLACEHOLDER: Describe the meteorological scenario — e.g., a marginal
%  severe-weather setup where parameters are near climatological
%  thresholds.]
The resulting possibility distribution is:
\[
\pi = \bigl\{\spc{MRGL}: 0.5,\; \spc{SLGT}: 0.5,\;
    \spc{ENH}: 0.0,\; \spc{MDT}: 0.0,\; \spc{HIGH}: 0.0\bigr\}
\]
with $\ign = 1 - 0.5 = 0.5$ and
$\Nc(\spc{MRGL}) = \Nc(\spc{SLGT}) = 0.0$.  Interpretation:
\spc{MRGL} and \spc{SLGT} are equally plausible ($\Pi = 0.5$ each),
high ignorance ($\ign = 0.5$) indicates that the system struggles with
this meteorological scenario, and neither outcome has conditional
necessity.

Blindly normalising Example~2 would inflate $\pi(\spc{MRGL})$ and
$\pi(\spc{SLGT})$ to 1.0, hiding the fact that the forecasting system
struggles with this meteorological scenario.  Preserving subnormality
makes these knowledge gaps explicit.

%----------------------------------------------------------
\subsection{Three-Component Uncertainty}\label{sec:three_component}

An analogue of ``classical'' necessity (requiring a normalised
possibility distribution $\pinorm$) supplements \poss{} and \ign{} for
each SPC category $A$ by computing necessity after normalisation of
known outcomes.  The three components are:

\begin{enumerate}[leftmargin=2em]
\item \textbf{Possibility, $\poss(A)$, $\pi$:} The raw (unnormalised)
    possibility of event $A$ from Equation~\eqref{eq:poss}, denoting
    evidence-based plausibility.
\item \textbf{Ignorance, \ign{}:} The global ignorance measure from
    Equation~\eqref{eq:ignorance}, indicating inherent deficiency of the
    forecasting system in capturing the notional ``true'' possibility
    distribution.
\item \textbf{Conditional Necessity, $\Nc(A)$, $\nu_c$:} Certainty of
    event $A$ computed \textit{after} normalising the distribution of
    known outcomes:
\begin{equation}
N_c(A) = 1 - \frac{\max_{\omega \notin A} \pi(\omega)}%
    {\max_{\omega \in \Omega} \pi(\omega)} \label{eq:cond_nec}
\end{equation}
This measures certainty conditional on the system's covered scenarios.
\end{enumerate}

All three quantities derive from the raw possibility distribution $\pi$:
\poss{} is the peak value for each category, \ign{} measures the global
gap to normality, and \Nc{} is computed after renormalisation.  Though
not statistically independent, they decompose the forecast signal into
complementary facets of decision-relevant information.

This framework bears structural resemblance to Dempster--Shafer theory's
belief--plausibility intervals: \poss{} corresponds to plausibility
(upper probability bound) and \Nc{} plays a role analogous to belief
(lower bound).  However, \ign{} has no direct Dempster--Shafer analogue,
as it quantifies model coverage gaps rather than evidential conflict
between sources.
% [PLACEHOLDER: Add Shafer (1976) citation.]

\textbf{Communication example (from Example~1):} ``\spc{MDT} risk is
the most plausible outcome ($\Pi = 0.75$) with moderate ignorance
($\ign = 0.25$).  Conditional necessity is $0.73$: among scenarios the
system covers, \spc{MDT} is the dominant category.''  This statement
conveys:
\begin{itemize}[leftmargin=2em]
\item \spc{MDT} risk is plausible but not guaranteed ($\Pi = 0.75$).
\item The system has moderate confidence in its assessment
    ($\ign = 0.25$).
\item Among scenarios the system \textit{does} cover, \spc{MDT} is the
    dominant outcome ($\Nc = 0.73$).
\end{itemize}

%----------------------------------------------------------
\subsection{Why Not Just Probabilities?}\label{sec:why_not_prob}

% [PLACEHOLDER: 3–4 paragraphs making the case for possibility theory
%  over pure probabilistic approaches.  Key arguments:
%  - Probabilities require a single number per outcome; possibility
%    distributions provide an interval [N(A), Pi(A)] that honestly
%    represents imprecision.
%  - The ignorance signal H_Pi has no probabilistic counterpart: a
%    probability forecast that sums to 1 cannot express "I don't know."
%  - In severe-weather forecasting, the SPC already uses categorical
%    language that maps more naturally onto possibility than probability:
%    "MDT risk is possible" vs "MDT risk has 40% probability."
%  - Possibility distributions can be converted to probabilities (via
%    the pignistic bridge, Section 3) but not vice versa — information
%    is lost in the probability-to-possibility direction.
%  - Practical benefit: a forecaster can issue a subnormal distribution
%    that explicitly communicates "today's setup is unusual and I am
%    less confident than normal" without having to manufacture precise
%    probability values.]

%----------------------------------------------------------
\subsection{Notation Table}\label{sec:notation_table}

Table~\ref{tab:notation} collects all recurring symbols used throughout
this paper.

\begin{table}[htbp]
\centering
\caption{Notation reference for possibilistic verification.}
\label{tab:notation}
\small
\begin{tabularx}{\columnwidth}{llX}
\toprule
Symbol & Macro & Meaning \\
\midrule
$\Omega$
    & ---
    & Universe of discourse (set of all outcomes) \\
$\pi$
    & ---
    & Raw (possibly subnormal) possibility distribution \\
$\pinorm$
    & \verb|\pinorm|
    & Normalised possibility distribution ($\pinorm = \pi / m$) \\
$m$
    & ---
    & Subnormality: $m = \max_{\omega \in \Omega} \pi(\omega)$ \\
$\poss(A)$
    & \verb|\poss|
    & Possibility of event $A$ (Eq.~\ref{eq:poss}) \\
$\nec(A)$
    & \verb|\nec|
    & Necessity of event $A$ (Eq.~\ref{eq:nec}) \\
$\Nc(A)$
    & \verb|\Nc|
    & Conditional necessity of $A$ (Eq.~\ref{eq:cond_nec}) \\
$\ign$
    & \verb|\ign|
    & Possibilistic ignorance (Eq.~\ref{eq:ignorance}) \\
$\alpha^*$
    & ---
    & Depth-of-truth: $\pinorm(c_{\mathrm{obs}})$ (Eq.~\ref{eq:alpha_star}) \\
$\eta$
    & ---
    & Nonspecificity (Eq.~\ref{eq:ns}) \\
$\delta$
    & ---
    & Resolution gap: $\alpha^* - \eta$ (Eq.~\ref{eq:rg}) \\
$\Nc^*$
    & ---
    & Conditional necessity of truth (Eq.~\ref{eq:nc_star}) \\
\bottomrule
\end{tabularx}
\end{table}
